%% ma: L11-B17-0019
%% lop: 11
%% chuong: 5
%% bai: 17
%% dang: 11.17.3
%% loai: TL
%% muc: VD
%% dap_an: "-"
%% nguon: "tai-lieu-hoc-tap-mon-toan-lop-11-theo-chuong-trinh-giao-duc-pho-thong-2018.pdf (Lớp Toán Thầy Duy, Phạm Lê Duy, THPT Chu Văn An), Chương 3 Bài 3, Dạng 3, Ví dụ 19 (tr.393-398)"
%% kien_thuc: "Lớp 11 Bài 17 (hàm số liên tục)"
%% trang_thai: chinh-thuc
%% ghi_chu: "Đã duyệt 03/11/2026."
\begin{ex}
Chứng minh rằng phương trình $a\cos2x+b\sin x+\cos x=0$ luôn có nghiệm với mọi tham số $a,b$.
\loigiai{\sloppy Đặt $f(x)=a\cos2x+b\sin x+\cos x$, liên tục trên $\mathbb R$. Ta có $f(0)=a+1$, $f\left(\dfrac\pi2\right)=-a+b$, $f(\pi)=a-1$, $f\left(\dfrac{3\pi}2\right)=-a-b$, và tổng bốn giá trị này bằng $0$. Nếu một giá trị bằng $0$ thì đó là nghiệm; nếu không thì trong bốn giá trị có số dương và số âm, tức có hai điểm $u,v\in\left[0;\dfrac{3\pi}2\right]$ với $f(u)f(v)<0$, nên phương trình có nghiệm giữa $u$ và $v$.}
\end{ex}
